\begin{abstract}
Language-model agents use experience through external retrieval and parameter learning. Consolidating stored experience into model weights requires deciding both which experience to train on and when its external records can leave retrieval. Admission based solely on student mastery can exclude useful teacher demonstrations, while removal immediately after training can discard guidance that the model still needs. We introduce \method, a framework that separates training admission from memory release. A demonstration-aware gate admits skills based on student mastery or sufficient successful teacher demonstrations. The student learns from selected action traces with context from an Engram Trace Graph (ETG). A verification gate then compares the updated model without episodic retrieval against the base model with retrieval on reserved probes. Release requires valid training participation and scores meeting both relative and absolute criteria; records for skills that fail verification remain retrievable. Experiments across interactive and reasoning benchmarks show consistent gains over frozen-weight ETG controls, with ablations supporting demonstration-aware admission and selective adapter activation. A controlled held-out comparison shows that verified release substantially reduces the retrievable footprint while remaining close to retain-all performance and outperforming immediate and budget-matched random release. \method\ also improves task performance while reducing input-token usage and model calls. Skill-level verification further reduces retrievable experience through selective archival, retaining external guidance for skills that do not meet the release criteria. Code is available in anonymous.4open.science/r/engram-anonymous.
\end{abstract}
